\section{Tuning a trajectory with vanishing perturbations}
\label{sec:trajectory}

We choose a temperature parameter for which the scale iteration remains
small and its potentials tend to zero in environmental moments.  This
is an existence argument for the iteration.  Its relation to the physical
critical point will be proved in Section~\ref{sec:criticality}.

\subsection{Initial couplings}

Write the dimensionless nearest-neighbor coupling as
\begin{equation}
 K(t)=\tfrac12\log(1+\sqrt2 e^t),\qquad
 K_e=K(h+\theta\xi_e),
 \qquad \theta>0.
 \label{traj:couplings}
\end{equation}
Thus $K(0)=K_0$ and
$\log((e^{2K_e}-1)/\sqrt2)=h+\theta\xi_e$.
The random variables $\xi_e$ are the independent symmetric signs of
the model.  This coordinate will also make duality particularly simple.

Fix an initial spacing $s_0$, large enough for the pure estimates and
the scale construction.  Allocate each bond to its nearest cell
center, splitting ties equally.  If $\alpha_{i,e}$ is the resulting
nonnegative allocation, then $\sum_i\alpha_{i,e}=1$, and the rule is
equivariant under grid translations and square symmetries.  Set
\begin{equation}
 f_i^{(0)}=\sum_e\alpha_{i,e}(K_e-K_0)b_e,
 \qquad b_e=\sigma_x\sigma_y-\E_0(\sigma_x\sigma_y).
 \label{traj:initial-potentials}
\end{equation}
These potentials are pure centered and their sum is precisely the
change in the Ising log weight, modulo a spin-independent constant.

\begin{lemma}[Initial moment bounds]\label{traj:initialization}
Let $a=K'(0)>0$.  For every fixed $T<\infty$, uniformly for
$|h|\le T\theta^2$, as $\theta\downarrow0$,
\begin{align}
 W_0&=2a^2\theta^2+O(\theta^3),&
 X_0&=O_{s_0}(\theta),\notag\\
 Y_0&=O_{s_0}(\theta^2+|h|),&
 \partial_hm_0(0,0)&=2s_0a>0.
 \label{traj:initial-bounds}
\end{align}
For every fixed cutoff $H>R_0$,
$v_{H,0}=2a^2\theta^2+O_H(\theta^3)$ on such an interval.
\end{lemma}
\begin{proof}
There are two bonds per lattice site.  Periodicity of the allocation
therefore gives $\sum_e\alpha_{0,e}=2s_0^2$.  Since $Q(b_e)=1$,
\[
 m_0=2s_0\E[K(h+\theta\xi)-K_0].
\]
Independence of distinct bond signs and $\sum_i\alpha_{i,e}=1$ give
the exact identity
\[
 W_0=s_0^{-2}\sum_e\alpha_{0,e}
       \Bigl(\sum_i\alpha_{i,e}\Bigr)\Var(K(h+\theta\xi))
     =2\Var(K(h+\theta\xi)).
\]
Taylor expansion of the smooth function $K$ proves the claimed
asymptotics and derivative.  The finite number of bounded bond
functions in a cell gives the bounds on $X_0$ and $Y_0$.
Lemma~\ref{var:near-W} gives the assertion about $v_{H,0}$.
\end{proof}

\subsection{The shooting argument}

The charge mean has an expanding linear term $Lm$, whereas the
variance has a negative quadratic correction.  We choose the initial
$h$ so that the expanding coordinate never leaves a narrow interval
of width proportional to the running variance.

\begin{theorem}[A small critical trajectory]
\label{traj:critical-trajectory}
Fix the scale-step parameters chosen in Proposition~\ref{mom:one-step}
and a sufficiently large initial spacing
$s_0$.  There are constants $A,B,D,c>0$, a fixed cutoff $H$, and
$\theta_0>0$ such that, for each $0<\theta<\theta_0$, one can choose
a deterministic $h=h(\theta)=O(\theta^2)$ with the following property.
The iterates
at spacings $s_k=s_0L^k$ satisfy, with $v_k=v_{H,k}$,
\begin{equation}
 v_k>0,\qquad
 X_k\le A\sqrt{v_k},\qquad
 Y_k\le Bv_k,\qquad
 |m_k|\le Dv_k,
 \label{traj:invariant-bounds}
\end{equation}
and
\begin{equation}
 v_k\le\frac{1}{v_0^{-1}+ck}.
 \label{traj:variance-decay}
\end{equation}
In particular $X_k,Y_k,m_k\to0$.  Their uniform bounds over all
$k$ can be made arbitrarily small by decreasing $\theta_0$.
\end{theorem}
\begin{proof}
We first choose constants so that the moment bounds propagate whenever
the mean remains between the walls $m=\pm Dv$.  In the estimates of
Proposition~\ref{mom:one-step}, write $C_1$ for a common constant
and retain its contraction constant $\rho<1$.  Choose $A$ large
enough for the initial $X_0/\sqrt{v_0}$ bound and so that
\begin{equation}
 \rho A+2C_1< A.
 \label{traj:A-choice}
\end{equation}
Choose $D$ next, with
\begin{equation}
 (L-1)D>C_1A^2+1.
 \label{traj:D-choice}
\end{equation}
Choose $T$ next so that, for every subsequently fixed $H$, the initial
ratios $m_0/(Dv_{H,0})$ at $h=-T\theta^2$ and $h=T\theta^2$ are
strictly below $-1$ and above $1$ for small $\theta$.
Lemma~\ref{traj:initialization} permits this choice independently
of $H$: the leading denominator is $2Da^2\theta^2$, whereas the
leading numerator divided by $\theta^2$ is an affine function of
$h/\theta^2$ with positive slope $2s_0a$.
Finally choose $B$ sufficiently large that
\begin{equation}
 \rho B+C_1(D+A^2)<B
 \label{traj:B-choice}
\end{equation}
and $Y_0\le Bv_{H,0}$ for $|h|\le T\theta^2$ and sufficiently small
$\theta$, for every subsequently fixed $H$.  This last choice is
possible from $Y_0\le C_{s_0}(1+T)\theta^2$ and
$v_{H,0}/\theta^2\to2a^2$.  The order so far is $A,D,T,B$.

For data satisfying \eqref{traj:invariant-bounds},
Lemma~\ref{var:near-W} gives
\begin{equation}
 W=v+O_H(v^{3/2}),
 \label{traj:W-comparison}
\end{equation}
where the implicit constant includes $A^3$.  Use
Proposition~\ref{var:flow} with $z=\sqrt v$.  Its constants $A,B$
are now fixed.  Choose $H$ sufficiently large that the $e_Hv^2$
error is smaller than one quarter of the negative leading term, using
the positive limit $2\pi\log L$ of the shell sum.  Then restrict $v$
to a sufficiently small interval, also within the smallness requirement
of Proposition~\ref{mom:one-step}, so that
\eqref{traj:W-comparison} and the $O_H(v^{5/2})$ remainder yield
\begin{equation}
 -C_2v^2\le v'-v\le-cv^2,
 \qquad v'>0,
 \qquad v'/v=1+O(v).
 \label{traj:scalar-recursion}
\end{equation}
Here $c,C_2>0$ are fixed and the smallness interval is uniform in
the scale.  Positivity follows, for example, by taking $C_2v<1/2$.

The $X$ estimate now gives
\[
 X'\le\rho A\sqrt v+C_1\sqrt W+C_1Bv
     \le\bigl(\rho A+C_1+o(1)\bigr)\sqrt v
     \le A\sqrt{v'}.
\]
The last inequality follows from \eqref{traj:A-choice} and
$v'/v\to1$.  Similarly,
\[
 Y'\le\bigl(\rho B+C_1(D+A^2)\bigr)v\le Bv'
\]
by \eqref{traj:B-choice}.  At the upper wall, the mean estimate gives
\[
 m'\ge LDv-C_1A^2v>Dv';
\]
at the lower wall it gives $m'<-Dv'$.  We decreased the smallness
interval if necessary to use \eqref{traj:D-choice} and
\eqref{traj:scalar-recursion}.  Thus the walls point outward, while
both norm bounds propagate up to and including the first exit step.

We now restrict $\theta$ so that the initial interval
$[-T\theta^2,T\theta^2]$ has all the properties fixed above and
lies in the uniform smallness range just obtained.  On this interval
$v_0=2a^2\theta^2+o(\theta^2)$ uniformly.  The initial $X$ bound is
covered by $A$, because changing $h$ by $O(\theta^2)$ does not change
the leading $O(\theta)$ bound on $X_0$.  Thus the complete parameter
order is $A,D,T,B,H$, followed by the small interval for $\theta$.

All moment data are continuous in $h$ at every finite number of
steps.  To check the only possible discontinuities, fix $h$ and
couple the constructions for nearby parameters with the same
environment, priorities, and threshold variables.  At a threshold
comparison, the threshold is independent of the data being compared
and has a continuous distribution.  Equality therefore has probability
zero.  At finitely many steps only finitely many comparisons enter
any specified local stencil, so the corresponding potentials are
almost surely continuous at the fixed $h$.  The crude real-log bounds
of Proposition~\ref{map:step} dominate them on every compact initial
parameter interval by a finite deterministic bound at each fixed
number of steps.  Dominated convergence applies to $X,Y,m,W$ and to
the finite polynomial $v_H$.

Continuity first supplies a closed initial subinterval $I_0$ on which
$-1\le m_0/(Dv_0)\le1$, with its two endpoints mapped to $-1$ and
$1$.  Suppose $I_k$ has been chosen so that all parameters in it
satisfy \eqref{traj:invariant-bounds} through time $k$, and the
endpoints lie on opposite time-$k$ walls.  The estimates just proved
give positive $v_{k+1}$ and the two norm bounds throughout $I_k$.
Its endpoints map strictly outside the opposite time-$(k+1)$ walls.
The intermediate value theorem therefore gives a closed subinterval
$I_{k+1}\subset I_k$ on which
$-1\le m_{k+1}/(Dv_{k+1})\le1$, with endpoints on the two walls.
For example, take a last crossing of the lower wall before a first
subsequent crossing of the upper wall.  The preceding-time bounds
remain valid on this smaller interval.  Compactness gives a parameter
$h\in\bigcap_{k\ge0}I_k$.  No monotonicity of an iterated coordinate
in $h$ is needed.

For this parameter, \eqref{traj:scalar-recursion} holds at every
step.  Taking reciprocals yields
\[
 \frac1{v_{k+1}}\ge\frac1{v_k}+c,
\]
and proves \eqref{traj:variance-decay}.  The other conclusions follow
from \eqref{traj:invariant-bounds} and $v_0=O(\theta^2)$.
\end{proof}

Every constant in Theorem~\ref{traj:critical-trajectory} is chosen
before a domain or a mesh approximation is specified.  The resulting
decay need not be summable over scales.  Section~\ref{sec:crossings}
will use the locality of each scale step to attach a geometric weight
to its contribution to a fixed macroscopic crossing test.
